\section{Related work}
\label{sec:related}

\subsection{BFT consensus protocols}

Classical Practical Byzantine Fault Tolerance~\cite{castro1999} tolerates
$f<\n/3$ faulty replicas at the cost of $O(\n^{2})$ messages per decision,
building on the impossibility and synchrony results
of~\cite{lamport1982,dwork1988}. HotStuff~\cite{yin2019} reduces message
complexity to $O(\n)$ through threshold signatures and pipelining;
Tendermint~\cite{buchman2016} provides instant finality;
Algorand~\cite{gilad2017} and HoneyBadgerBFT~\cite{miller2016} employ
randomization; Mir-BFT~\cite{stathakopoulou2019} parallelizes leaders;
Ouroboros~\cite{kiayias2017} gives provable security for stake-based selection.
Vukoli\'{c}~\cite{vukolic2015} argues that no single protocol dominates across
operating regimes. All of these works fix the validator set at deployment time
and none provides a procedure for choosing its size.

\subsection{Blockchain performance benchmarking}

Empirical characterizations of permissioned ledgers include Hyperledger
Fabric~\cite{androulaki2018,thakkar2018}, Quorum~\cite{baliga2018}, the
BLOCKBENCH framework~\cite{dinh2017}, and the BFT-SMaRt ordering
service~\cite{sousa2018}. These studies report throughput at a handful of fixed
configurations. They do not separate client-side offered load from validator
count, which is exactly the confounding we formalize in
Theorem~\ref{thm:identifiability}, and they do not extrapolate.

\subsection{Scaling laws for parallel and distributed systems}

Amdahl's law~\cite{amdahl1967} and Gustafson's reformulation~\cite{gustafson1988}
bound speedup under a fixed serial fraction. Gunther's Universal Scalability
Law~\cite{gunther2007,gunther2015} adds a coherency-delay term quadratic in the
number of participants,
\begin{equation}
C(\n)=\frac{\n}{1+\sigma(\n-1)+\kappa\,\n(\n-1)},
\label{eq:usl}
\end{equation}
and is the closest existing model to ours. USL describes \emph{capacity as a
function of the number of workers sharing one workload}; it does not separate
the offered-concurrency dimension, and it has not been calibrated for BFT
consensus. We use~\eqref{eq:usl} as a baseline in
Section~\ref{sec:results}.

\subsection{Concentration bounds and statistical detection}

Hoeffding's inequality~\cite{hoeffding1963} and the broader concentration
toolkit~\cite{boucheron2013} give exponential tail bounds for bounded
independent averages, which we use to bound the probability that a faulty
minority escapes detection. Prediction intervals for nonlinear regression follow
the delta method~\cite{seber2003,bates1988}; the residual bootstrap
follows~\cite{efron1993,davison1997}.

\subsection{Optimization under uncertainty}

Chance-constrained programming originates with Charnes and
Cooper~\cite{charnes1959}; tractable convex approximations are surveyed
in~\cite{nemirovski2007,bental2009}. To our knowledge, chance constraints have
not previously been applied to validator-set sizing, where the novelty lies in
the \emph{opposing monotonicity} of the performance and safety constraints,
which is what yields the closed feasibility interval of
Theorem~\ref{thm:window}.

\subsection{Blockchain in e-governance}

Security analyses of deployed remote-voting systems~\cite{springall2014,%
specter2020} and the critical assessment of blockchain voting
in~\cite{park2021} establish that consensus-layer guarantees do not resolve
endpoint or coercion threats. Privacy techniques such
as~\cite{bensasson2014} are complementary. Our earlier
work~\cite{peniak2026a,peniak2026b} developed, respectively, a single-factor
throughput calibration procedure and an adaptive consensus-switching mechanism
with detection-based safety bounds. The present paper unifies and corrects
these two lines: it identifies the confounding that inflated the single-factor
exponent of~\cite{peniak2026a}, and it supplies the empirical detector
characterization that~\cite{peniak2026b} assumed.

\subsection{Positioning}

Table~\ref{tab:related} summarizes the comparison.

\begin{table}[!t]
\caption{Positioning with respect to representative prior work}
\label{tab:related}
\centering
\small
\begin{tabular}{lccccc}
\headline
Approach & Scaling & Client & Prediction & Safety & Sizing \\
         & model   & factor & intervals  & bound  & rule   \\
\headline
PBFT~\cite{castro1999}              & --      & --  & --  & yes & -- \\
HotStuff~\cite{yin2019}             & $O(\n)$ & --  & --  & yes & -- \\
BLOCKBENCH~\cite{dinh2017}          & --      & --  & --  & --  & -- \\
Fabric bench.~\cite{thakkar2018}    & --      & partial & -- & -- & -- \\
USL~\cite{gunther2007}              & yes     & --  & --  & --  & -- \\
Single-factor~\cite{peniak2026a}    & yes     & --  & heuristic & -- & -- \\
Adaptive consensus~\cite{peniak2026b} & --    & --  & --  & yes & -- \\
\headline
This work                           & yes     & yes & calibrated & yes & yes \\
\headline
\end{tabular}
\end{table}
